\section{Estado del arte}

\subsection{\normalsize Simuladores cuánticos}
La simulación clásica de circuitos cuánticos es esencial para el desarrollo y
validación de algoritmos cuánticos. Sin embargo, el principal desafío radica
en la representación del estado cuántico, cuyo tamaño crece exponencialmente con
el número de qubits.

Uno de los primeros enfoques para abordar este problema fue la representación
mediante \textit{Matrix Product States} (MPS), propuesta por \textcite{Vidal2003}.
Este método permite representar ciertos estados cuánticos de forma eficiente
cuando el entrelazamiento es limitado, reduciendo los requisitos de almacenamiento
y cómputo. Su principal ventaja es la escalabilidad en circuitos con baja
profundidad, permitiendo simular cientos de qubits. No obstante, su limitación
radica en que para sistemas altamente entrelazados el costo computacional crece
exponencialmente, restringiendo su aplicabilidad a ciertos algoritmos.

A medida que se buscaban métodos más generales, surgió la simulación tipo
\textit{Schrödinger}, que almacena el vector de estado completo y lo actualiza
en cada operación cuántica. Este enfoque, empleado en simuladores como QuEST y
Qiskit Aer \cite{Jones2019, qiskit_aer_doc}, garantiza alta precisión y permite simular cualquier
circuito cuántico. Su desventaja principal es la extrema demanda de memoria, la
cual crece exponencialmente con el número de qubits: una simulación de 45 qubits
ya ha requerido 0.5 petabytes de memoria \cite{Haner2017}, y al extender esta
representación de doble precisión a 50 qubits el costo pasa al orden de decenas de
petabytes.

Para mitigar el problema de memoria, se desarrollaron métodos basados en
partición del circuito y caminos de Feynman. Por ejemplo, Chen y cols.\
presentaron un algoritmo distribuido capaz
de calcular amplitudes de salida y simular circuitos universales de hasta
$12\times12$ qubits con profundidad moderada sin materializar el estado completo
\cite{Chen2018}. Sin embargo, aunque reducen el
almacenamiento requerido, estos métodos incrementan la complejidad computacional,
ya que el número de caminos crece exponencialmente con la profundidad del
circuito.

\subsection{\normalsize Compresión de memoria con pérdida de precisión}
Dado que la representación exacta de estados cuánticos es ineficiente en términos
de memoria, se han explorado técnicas de compresión con pérdida de precisión. Estos
enfoques sacrifican una pequeña cantidad de fidelidad numérica a cambio de una
reducción significativa en el uso de memoria, permitiendo la simulación de
circuitos más grandes.

Uno de los primeros enfoques en este ámbito fue el de \textcite{wu2018amplitude, wu2019full},
quienes introdujeron la compresión adaptativa con error acotado. Utilizando la
biblioteca SZ, lograron reducir el tamaño del vector de estado hasta en un factor
de 16, manteniendo fidelidades superiores al 99.9~\%. Su principal ventaja es
que permite ampliar el número de qubits simulados sin un incremento
proporcional en los requisitos de memoria. Sin embargo, la compresión y
descompresión recurrente introduce sobrecarga computacional, afectando la eficiencia
temporal de la simulación.

En un esfuerzo por maximizar la escala de la simulación cuántica de estado
completo, \textcite{zhang2024bmqsim} presentaron \emph{BMQSim}, un marco que
extiende \emph{SV-Sim} mediante el uso de \textbf{compresores con pérdida de
precisión y control de error punto a punto} ejecutados directamente en GPU (por
ejemplo, variantes de la familia SZ/cuSZ)\cite{zhang2024bmqsim}. Su diseño
combina:
\begin{itemize}
    \item Una estrategia de \textit{circuit partitioning} que disminuye la
        frecuencia de (des)compresión.

    \item Un \textit{pipeline} solapado que oculta en gran medida el coste de
        mover y comprimir datos.

    \item Una gestión de memoria jerárquica (VRAM + SSD mediante GPUDirect~Storage)
        que asigna dinámicamente espacio a los bloques comprimidos.
\end{itemize}
Con estas técnicas, BMQSim reduce el consumo de memoria en más de un orden de magnitud
y mantiene fidelidades superiores al 99~\% sin penalizar de forma apreciable el
tiempo de simulación. Su eficacia, sin embargo, está condicionada a la
disponibilidad de hardware acelerado con GPUs de alto rendimiento y
unidades NVMe rápidas, lo que puede restringir su adopción en plataformas de menor
capacidad computacional.

Más recientemente, \textcite{Huffman2024} exploraron la cuantización de
amplitudes como un mecanismo para reducir la precisión numérica sin afectar la
fidelidad global del sistema. Demostraron que representaciones con tan solo 7 bits
por amplitud permiten mantener una fidelidad superior al 99~\%, lo que sugiere que
la precisión completa de punto flotante puede ser innecesaria en muchas simulaciones.
Aunque esta técnica ofrece un enfoque simple y eficiente para reducir el uso
de memoria, su aplicabilidad a circuitos de gran escala aún requiere validación
empírica.

\textcite{DiazToro2025} propone un esquema de \textbf{vector de estado
comprimido} que emplea compresores con pérdida de precisión—principalmente
\emph{ZFP} en modo \emph{fixed‑rate}, complementado por \emph{SZ}—para reducir
la huella de memoria del simulador. El autor reporta reducciones del \textbf{10--20~\%}
en casos base y, en configuraciones de mayor escala, ahorros que alcanzan hasta
un \textbf{75~\%} de la memoria requerida. \emph{ZFP} opera sobre bloques independientes
de $4^{d}$ valores, de modo que la (des)compresión de cada bloque se realiza en
completo aislamiento del resto; esto permite solapar dichas operaciones con el
cómputo principal y amortiguar la sobrecarga introducida. Aunque la compresión
incrementa el tiempo de ejecución, el impacto se compensa al transmitir los
fragmentos del vector en formato comprimido dentro de una arquitectura
distribuida basada en \emph{MPI}, reduciendo el volumen de comunicación y
habilitando simulaciones de hasta $30$ qubits en los experimentos reportados. Finalmente,
el estudio concluye que esta estrategia de \textit{estado completo comprimido
distribuido} resulta más fiable y escalable que la \textit{poda de estados de
baja probabilidad}, la cual presenta descensos de fidelidad y sobrecarga adicional
que la hacen desaconsejable.

\subsection{\normalsize{Compresión de memoria sin pérdida de precisión}}
La simulación cuántica eficiente requiere reducir el consumo de memoria sin
comprometer la fidelidad del estado cuántico. Para ello, se han desarrollado
técnicas de compresión sin pérdida de precisión que buscan representar estados
y operadores cuánticos de manera más compacta sin alterar sus valores.

Uno de los primeros enfoques en esta dirección fue el uso de diagramas de decisión
cuánticos (\textit{Quantum Information Decision Diagrams}, QuIDD),
introducidos por \textcite{Viamontes2003}. Este método permite representar
vectores de estado y operadores unitarios aprovechando redundancias
estructurales en sus coeficientes, lo que permite una representación más eficiente
en memoria. Su ventaja radica en que, para circuitos con alta redundancia, el
crecimiento en memoria puede reducirse de exponencial a polinomial. Sin embargo,
su eficacia se limita a circuitos con estructura repetitiva, fallando en casos
de entrelazamiento complejo.

En un desarrollo posterior, \textcite{Miller2006} propusieron los \textit{Quantum
Multiple-valued Decision Diagrams} (QMDD), los cuales extienden la
representación de QuIDD mediante el uso de aristas ponderadas con matrices
unitarias. Esto permitió manejar de manera más eficiente circuitos reversibles
y operadores cuánticos de mayor tamaño. Aunque estos diagramas optimizan la
representación de puertas lógicas y matrices densas, su eficiencia sigue dependiendo
de la estructura interna del circuito.

Más recientemente, \textcite{Jaques2021} exploraron la esparsidad del estado
cuántico para almacenar solo las amplitudes no nulas, reduciendo significativamente
el almacenamiento requerido en algoritmos con exploración limitada del espacio
de Hilbert. Su enfoque demostró ser altamente eficiente en problemas de
factorización y aritmética cuántica, pero pierde efectividad en circuitos con alta
interferencia y entrelazamiento.

En 2023, \textcite{Vinkhuijzen2023} introdujeron los \textit{Local Invertible
Map Decision Diagrams} (LIMDD), una evolución de los diagramas de decisión que
incorpora transformaciones locales para mejorar la representación compacta del
estado cuántico. Esta técnica permite manejar estados estabilizadores y ciertas
estructuras altamente entrelazadas que eran difíciles de representar con diagramas
de decisión tradicionales. No obstante, su aplicabilidad sigue limitada a casos
donde las transformaciones locales puedan ser explotadas para reducir la complejidad
estructural.

Estos avances en compresión sin pérdida han permitido extender la capacidad de
simulación de circuitos cuánticos en hardware clásico. Sin embargo, aún
existen limitaciones en la representación de estados arbitrarios, lo que ha llevado
a la exploración de otros métodos para optimizar el uso de memoria.

\subsection{\normalsize{Enfoques de compresión sin pérdida para datos de punto flotante}}

En simulación científica es frecuente trabajar con arreglos extensos de números de
punto flotante, cuyo volumen puede convertir el almacenamiento y el movimiento de
datos en una restricción de desempeño. En este contexto, la compresión sin pérdida
busca reducir el tamaño de la representación binaria sin alterar ningún bit de la
información original, de modo que la reconstrucción coincida exactamente con los datos
de entrada. Esta exigencia es especialmente relevante en escenarios donde los valores
comprimidos participan posteriormente en nuevas etapas de cálculo o sirven como
referencia numérica exacta.

A diferencia de la compresión de texto o de datos discretos, la compresión de punto
flotante presenta dificultades particulares. Muchos arreglos científicos contienen
redundancias asociadas a continuidad espacial, dependencia temporal o regularidades en
la representación binaria de signo, exponente y mantisa. Sin embargo, cuando esas
regularidades son débiles o no siguen una estructura local simple, la compresión se
vuelve más desafiante. Por esta razón, la literatura ha propuesto distintos enfoques que
no actúan todos del mismo modo: algunos se basan en predicción, otros en
transformaciones reversibles de la representación binaria y otros en reorganización
tipada de los datos antes de aplicar un códec general.

\subsubsection{FPC (Burtscher \& Ratanaworabhan)}

\textit{FPC} (\emph{Floating Point Compressor})\cite{burtscher2009fpc} es un algoritmo
de compresión sin pérdida diseñado para secuencias lineales de valores de punto
flotante, especialmente en doble precisión. Su principio central consiste en predecir el
siguiente valor a partir del historial reciente y comprimir únicamente la diferencia entre el
valor real y el valor estimado.

Para ello, FPC emplea dos predictores derivados de esquemas tipo FCM/DFCM, cuyos
resultados se comparan con el dato real mediante una operación \texttt{XOR}. Cuando la
predicción es cercana, la diferencia resultante presenta una cantidad considerable de
bytes iniciales nulos, que pueden codificarse de forma compacta. El algoritmo almacena
entonces un pequeño código de control junto con los bytes residuales no nulos. De esta
manera, la compresión depende directamente de la capacidad del predictor para anticipar
la estructura local de la secuencia.

Este enfoque resulta especialmente adecuado cuando existe correlación entre valores
adyacentes o cuando la evolución del arreglo mantiene cierta regularidad. En cambio, su
efectividad disminuye cuando los datos no exhiben continuidad local suficiente para que
la predicción produzca residuos pequeños. Aun así, FPC constituye una referencia
importante en compresión sin pérdida de punto flotante porque muestra con claridad la
utilidad y las limitaciones de los esquemas predictivos sobre representaciones IEEE 754.

\subsubsection{FPZIP (Lindstrom \& Isenburg)}

\textit{FPZIP}\cite{lindstrom2006fpzip} es un compresor sin pérdida orientado a arreglos
de punto flotante en contextos científicos. Su diseño parte de la idea de que muchos
conjuntos de datos numéricos presentan continuidad o dependencia local, de modo que
pueden describirse con mayor eficiencia a partir de errores de predicción que a partir de
los valores originales.

A diferencia de enfoques lineales simples, FPZIP incorpora esquemas de predicción
adaptados a la estructura de los datos, lo que le permite aplicarse a arreglos de
distintas dimensionalidades. Tras generar una estimación para cada valor, el algoritmo
codifica el residuo de predicción mediante un modelo entrópico diseñado para preservar
exactitud bit a bit sobre la representación en punto flotante. En este sentido, FPZIP no
requiere cuantización ni aproximación intermedia: la compresión sigue siendo
estrictamente reversible.

La relevancia de FPZIP radica en que explota de manera más
general la estructura local de arreglos científicos y combina predicción con codificación
entrópica especializada. Esto lo convierte en un referente importante para datos en los
que la redundancia no se manifiesta únicamente entre valores consecutivos, sino también
en relaciones espaciales o estructurales más amplias.

\subsubsection{Bitshuffle con LZ4 (Masui et al.)}

Una estrategia diferente consiste en no predecir los valores, sino transformar su
disposición binaria para hacer visibles regularidades que no son evidentes en el orden
original. Bajo esta idea, \textit{Bitshuffle}\cite{masui2015bitshuffle} actúa como un filtro
reversible que reorganiza los bits de una secuencia de valores de punto flotante antes de
aplicar un códec de compresión general, usualmente LZ4.

El procedimiento puede entenderse como una transposición bit a bit: si se consideran
$N$ valores de $m$ bits, sus representaciones binarias forman una matriz de
$N \times m$, y Bitshuffle reordena los datos agrupando los bits según su posición. Esta
transformación puede revelar repeticiones o patrones en determinados planos de bits,
por ejemplo en exponentes o signos, que de otra forma quedarían dispersos entre los
bytes de cada valor. Una vez realizada la transposición, un códec LZ puede aprovechar
con mayor eficacia esas secuencias más homogéneas.

Este enfoque no depende de predicciones locales entre valores adyacentes, sino de la
posibilidad de exponer regularidades internas de la representación binaria. Por ello, su
interés en compresión científica es doble: por un lado, preserva completamente la
exactitud; por otro, ofrece una vía distinta para explotar redundancia en datos donde los
modelos predictivos no siempre resultan adecuados. En la práctica, Bitshuffle suele
entenderse mejor como una transformación previa al códec que como un compresor
autosuficiente.

\subsubsection{Transformación de Datos Tipados (TDT)}

La \textit{Transformación de Datos Tipados} (TDT)\cite{jamalidinan2025tdt} representa
una línea más reciente de trabajo orientada a reorganizar datos de punto flotante antes
de aplicar compresión generalista. A diferencia de Bitshuffle, que opera sobre una
transposición fija de bits, TDT busca identificar posiciones de byte con comportamientos
estadísticos diferentes y reagruparlas para producir flujos más homogéneos.

La idea central es que, en una representación IEEE 754, no todos los bytes cumplen el
mismo papel ni presentan el mismo nivel de entropía. Los bytes asociados al exponente,
al signo o a distintas regiones de la mantisa pueden exhibir distribuciones distintas
según la naturaleza de los datos. TDT analiza esas diferencias a nivel de bloque y aplica
una reorganización reversible basada en características estadísticas, de modo que los
bytes con comportamiento semejante queden contiguos antes de ser entregados a un
compresor convencional.

Más que un códec independiente, TDT debe entenderse como un preprocesamiento
tipado orientado a mejorar la eficiencia de compresores generales sobre datos
numéricos. Su relevancia dentro del estado del arte radica en mostrar que la compresión
sin pérdida de flotantes no depende únicamente de predicción o de compresión
entrópica, sino también de la capacidad de reordenar la representación binaria para
hacer más visible su estructura estadística.

\subsection{\normalsize{Otros métodos de optimización de memoria}}
Además de la compresión sin pérdida, se han desarrollado enfoques alternativos
para reducir el consumo de memoria en simuladores cuánticos, incluyendo
técnicas basadas en redes tensoriales, poda de estados y optimización HPC.

Uno de los primeros enfoques en esta dirección fue propuesto por \textcite{Markov2008},
quienes introdujeron el uso de redes tensoriales para simular circuitos
cuánticos mediante contracción optimizada de tensores. En este método, el
circuito se representa como un grafo tensorial, y su simulación se optimiza contrayendo
secuencias de tensores en un orden eficiente. Esto permite representar ciertos
circuitos con memoria subexponencial, aunque la efectividad del método depende
de la estructura del entrelazamiento.

Posteriormente, \textcite{Boixo2017} propusieron un modelo gráfico para
circuitos cuánticos de baja profundidad basado en la eliminación de variables
en grafos probabilísticos. Su técnica permitió simular circuitos aleatorios cercanos
al régimen de supremacía cuántica en hardware clásico, extendiendo el tamaño
de circuitos simulables sin necesidad de almacenar el estado cuántico completo.
Sin embargo, este método pierde eficacia conforme aumenta la profundidad del
circuito.

En paralelo, \textcite{Pednault2017} introdujeron técnicas de diferimiento de contracciones
tensoriales en simulaciones HPC. Su propuesta permite manejar circuitos más
grandes al intercambiar memoria por cómputo adicional, almacenando resultados intermedios
en memoria secundaria y evitando la expansión total del estado cuántico en RAM.
Esta estrategia ha sido clave en la simulación de circuitos con estructuras
altamente no triviales, aunque su implementación requiere acceso a infraestructuras
de supercomputación.

Finalmente, como se analizó previamente, la tesis doctoral de \textcite{DiazToro2025}
introduce técnicas de compresión con pérdida de precisión en esquemas distribuidos
de simulación cuántica. Este antecedente resulta especialmente relevante porque
consolida el uso de bibliotecas científicas de compresión, vectores de estado y
paralelismo distribuido con \emph{MPI} como una línea válida para reducir consumo
de memoria en simulación cuántica. Al mismo tiempo, deja visible un espacio de
estudio complementario: evaluar hasta qué punto técnicas de compresión
\textbf{sin pérdida de precisión} pueden aportar reducciones útiles de memoria sin
introducir discrepancias numéricas en el estado simulado.
